\section{Empirical Results}
\label{sec:results}

\subsection{Experimental Setup}
\label{sec:results:setup}

We evaluate BisectionEnsemble on three states chosen to represent qualitatively
distinct redistricting problems:
\begin{itemize}
  \item \textbf{North Carolina} ($k=14 = 7 \times 2$): a competitive swing state
        with a mixed binary/prime factorization. BisectionEnsemble applies at the
        7 bisection nodes at depth 2.
  \item \textbf{Wisconsin} ($k=8 = 2^3$): a fully binary tree of depth 3.
        BisectionEnsemble applies at all 7 tree nodes (1 root + 2 at depth 1 + 4 at depth 2).
  \item \textbf{Texas} ($k=38 = 2 \times 19$): a large state where direct
        GerryChain $k$-way ensemble fails. The AR factorization is $19 \times 2$
        (19-way primary METIS split, then bisect each of the 19 sub-regions).
        BisectionEnsemble applies at the 19 bisection nodes at depth 2.
\end{itemize}

For each state, we compare three methods:
\begin{enumerate}
  \item \textbf{Standard bisection}: Single METIS call at each tree node.
        This is the baseline from ApportionRegions \citep{dellaLibera2026apportion}.
  \item \textbf{BisectionEnsemble($p$=0.5, $T$=100)}: Local 100-step ReCom at
        each 2-way node, selecting the median edge-cut plan.
  \item \textbf{BisectionEnsemble($p$=0.0, $T$=100)}: Local 100-step ReCom at
        each 2-way node, selecting the minimum edge-cut plan.
\end{enumerate}
We also attempt direct GerryChain $k$-way ensemble on all three states for
comparison, recording success/failure and (for successful states) acceptance rate.

All experiments use the 2020 Census P.L.~94-171 redistricting data.
Partisan outcomes are measured against 2020 presidential two-party precinct
returns from the Voting and Election Science Team (VEST) Harvard Dataverse
\citep{vest2020}, interpolated to census tracts via areal weighting.
Polsby-Popper compactness is computed using the standard formula
$\text{PP} = 4\pi A / P^2$ where $A$ is district area and $P$ is perimeter
\citep{polsby1991third}.
Runtimes are measured on a 16-core workstation with Rayon parallelism enabled.

\subsection{Bipartition Failure: GerryChain vs.\ BisectionEnsemble}
\label{sec:results:failure}

Table~\ref{tab:failure} records the bipartition failure rate for direct
GerryChain $k$-way ensemble attempts at $T=100$ steps per state.

\begin{table}[h]
\centering
\begin{tabular}{lrrrr}
\toprule
State & $k$ & GerryChain acceptance rate & GerryChain status & BisectionEnsemble status \\
\midrule
North Carolina & 14 & 61\% & Success & Success \\
Wisconsin      & 8  & 84\% & Success & Success \\
Texas          & 38 & $<$1\% & Failure (stalled) & Success \\
\bottomrule
\end{tabular}
\caption{Bipartition failure comparison. GerryChain stalls on Texas ($k=38$):
  fewer than 1\% of ReCom steps are accepted at 100 steps, yielding 0--1
  accepted plans and no usable distribution. BisectionEnsemble succeeds on
  all three states because every local ensemble is 2-way.}
\label{tab:failure}
\end{table}

GerryChain succeeds on NC ($k=14$) and WI ($k=8$) but stalls on TX ($k=38$)
with an acceptance rate below 1\%. The correct causal account (described in
Section~\ref{sec:intro}) is geometric, not a balance-ratio issue: GerryChain's
ReCom step merges two adjacent districts and re-bipartitions the merged region
into two equal halves. For $k=38$, the merged region contains $\approx 2/38$
of the state population, but the critical difficulty is that the spanning tree
of the merged region must be cut to achieve a 1:1 split within that region,
and the geometric constraints from the 36 surrounding districts severely limit
the set of valid spanning tree cut edges. Balanced spanning tree cuts of a
geometrically constrained region are rare, producing near-zero acceptance.
\emph{Note: this is not a 1:37 population ratio problem --- the 1:1 balance is
within the merged pair, but valid cut edges are sparse due to the boundary
geometry imposed by the surrounding districts.}
BisectionEnsemble avoids this entirely: every local ReCom
step splits a subregion into two roughly equal halves at a bisection node,
and the absence of surrounding-district boundary constraints (the subregion is
treated as an isolated local graph) makes balanced cuts abundant.

\emph{Reproducibility note (Phase 2 deferred):} The GerryChain $<$1\%
acceptance figure for TX is a single-run observation. GerryChain experiment
reproducibility details --- version, configuration, number of runs, hardware,
and random seed --- are deferred to the Phase~2 empirical study. Phase~2 will
report the distribution of acceptance rates across 10 independent runs to
establish whether the near-zero acceptance is stable or run-dependent.

\subsection{Compactness Comparison}
\label{sec:results:compactness}

Table~\ref{tab:compactness} reports mean Polsby-Popper compactness across
all districts in each state's plan.

\begin{table}[h]
\centering
\begin{tabular}{lrrrrr}
\toprule
State & $k$ & Standard bisection & BE($p$=0.5) & BE($p$=0.0) & GerryChain \\
\midrule
North Carolina & 14 & 0.217 & 0.242 (+11.5\%) & 0.261 (+20.3\%) & 0.249 \\
Wisconsin      & 8  & 0.198 & 0.209 (+5.6\%)  & 0.228 (+15.2\%) & 0.221 \\
Texas          & 38 & 0.184 & 0.206 (+11.9\%) & 0.231 (+25.5\%) & --- (failed) \\
\bottomrule
\end{tabular}
\caption{Mean Polsby-Popper compactness. BE($p$=0.5) denotes BisectionEnsemble
  at median percentile; BE($p$=0.0) at minimum edge-cut.
  Standard bisection is the baseline. GerryChain entry for TX is unavailable
  due to bipartition failure. Parenthetical percentages are improvements over
  standard bisection.}
\label{tab:compactness}
\end{table}

BisectionEnsemble consistently improves compactness over standard METIS bisection.
The median-percentile variant ($p=0.5$) provides 5--12\% improvement; the
minimum-cut variant ($p=0.0$) provides 15--25\% improvement. These gains arise
because the local ReCom chain explores bisections that METIS, with its
multilevel coarsening heuristic, may not reach in a single run.

The minimum-cut variant ($p=0.0$) achieves the best compactness by definition:
it selects the bisection with the smallest boundary length among all accepted
plans, which directly minimises district perimeter and thus maximises
Polsby-Popper scores. The median variant ($p=0.5$) is more representative of
the feasible bisection space and may be preferable when compactness optimisation
is not the primary objective.

\subsection{Partisan Outcomes}
\label{sec:results:partisan}

Table~\ref{tab:partisan} reports partisan outcomes for each method on each state.

\begin{table}[h]
\centering
\begin{tabular}{lrrrr}
\toprule
State & $k$ & Standard bisection & BE($p$=0.5) & BE($p$=0.0) \\
\midrule
North Carolina & 14 & 7D/7R & 7D/7R & 7D/7R \\
Wisconsin      & 8  & 2D/6R & 3D/5R & 2D/6R \\
Texas          & 38 & 16D/22R & 16D/22R & 15D/23R \\
\bottomrule
\end{tabular}
\caption{Partisan outcomes (2020 presidential returns). All outcomes are
  Democratic seats / Republican seats. NC outcomes are identical across all
  methods. WI shows a partisan shift at $p=0.5$ but not $p=0.0$, reflecting
  that median-cut bisections find a different equilibrium than the single METIS
  bisection. TX is stable between standard and BE($p=0.5$) but shifts by one seat
  at $p=0.0$.}
\label{tab:partisan}
\end{table}

The key finding is that BisectionEnsemble does not systematically shift
partisan outcomes in any direction. NC remains 7D/7R across all methods,
reflecting the geographic robustness of its near-proportional outcome
established in \citet{dellaLibera2026apportion}. WI shows a one-seat shift
at $p=0.5$: the median-cut bisection at the three root-level nodes produces
a slightly different spatial arrangement of competitive tracts, pushing one
district from safe-Republican to lean-Democratic. This shift disappears at
$p=0.0$, where the minimum-cut bisection returns to the same local optimum
as standard METIS.

\subsection{Runtime}
\label{sec:results:runtime}

Table~\ref{tab:runtime} reports wall-clock runtime for each method with
Rayon parallelism across 16 cores.

\begin{table}[h]
\centering
\begin{tabular}{lrrrr}
\toprule
State & $k$ & Standard bisection & BE($p$=0.5, $T$=100) & BE($p$=0.5, $T$=500) \\
\midrule
North Carolina & 14 & 0.8\,s & 4.3\,s & 20.1\,s \\
Wisconsin      & 8  & 0.4\,s & 3.1\,s & 14.8\,s \\
Texas          & 38 & 2.1\,s & 9.7\,s & 47.2\,s \\
\bottomrule
\end{tabular}
\caption{Wall-clock runtime (16 cores, Rayon parallelism). Standard bisection
  is a single METIS call per node. BE($T$=100) runs 100 ReCom steps per node
  with full parallelism across same-depth nodes. The $T$-fold runtime scaling
  is approximate due to parallelism and load balancing effects.}
\label{tab:runtime}
\end{table}

BisectionEnsemble at $T=100$ steps takes 4--10$\times$ longer than standard
bisection. This is expected: $T=100$ Wilson's algorithm runs per node add
$O(T \cdot n_\text{node})$ work per node. The runtime is practical for
operational use: the full BisectionEnsemble run for Texas ($k=38$, 5{,}200
tracts, $T=100$) completes in under 10 seconds on 16 cores, compared to
standard bisection's 2.1 seconds. For $T=500$ (which provides a broader
ensemble distribution), runtime remains under 1 minute for all three states.

\paragraph{Parallelism efficiency.}
The 19 bisection nodes at depth 2 in the Texas tree (each managing $\approx 140$
tracts) run fully in parallel after the depth-1 METIS call completes.
With 16 cores and 19 nodes, load balance is near-perfect (19/16 $\approx$ 1.2
imbalance ratio). The depth-1 METIS call (19-way partition of the full state)
is not parallelised and takes 0.3\,s, setting the sequential lower bound.

\paragraph{Phase 2 empirical study (deferred).}
Tables~\ref{tab:compactness}, \ref{tab:partisan}, and \ref{tab:runtime} present
point estimates from single BisectionEnsemble runs per method per state.
For a stochastic algorithm, single-run estimates are insufficient to characterise
the method's distribution of outcomes. Phase~2 will extend these tables to
multi-run statistics: mean $\pm$ standard deviation across 10 independent runs,
each using a different CSPRNG-drawn METIS initialisation seed. In particular,
the WI one-seat shift at $p=0.5$ (Table~\ref{tab:partisan}) should be interpreted
as a single-run observation pending multi-run confirmation of its stability.
A fourth baseline, MultiSeedMETIS($N=100$) --- which runs METIS 100 times with
different seeds at each node and selects the plan at rank $\lfloor p \times 100
\rfloor$ --- will also be included in Phase~2 to isolate the contribution of
the ReCom chain from the contribution of multiple-initialisation sampling.
Additionally, Phase~2 will extend the empirical evaluation to at least 5--7
states (including strongly partisan and additional prime-$k$ states) to establish
whether the compactness and partisan-neutrality findings generalise.
